% Frequency-aware diffusion losses: what a spectral correction is worth.
%
% Deliberately plain `article` rather than a venue template. The target venue
% is not settled, and swapping \documentclass late is far cheaper than
% unpicking a template's own section and float conventions from the prose.
% To move to LNCS: change the class, replace \maketitle's block, keep sections.
%
% Build:  bash paper/build.sh      (or: latexmk -pdf main.tex)

\documentclass[11pt,a4paper]{article}

\usepackage[utf8]{inputenc}
\usepackage[T1]{fontenc}
\usepackage[margin=2.5cm]{geometry}
\usepackage{amsmath,amssymb}
\usepackage{graphicx}
\usepackage{booktabs}
\usepackage{siunitx}
\usepackage{xcolor}
\usepackage[numbers,sort&compress]{natbib}
\usepackage[hidelinks]{hyperref}

\graphicspath{{figures/}}

% Numbers that appear in more than one place. Defining them once means a
% re-run updates every mention, rather than leaving one stale figure in the
% abstract to be found by a referee.
\newcommand{\deficitHigh}{2.20}          % MSE high-band deficit, dB
\newcommand{\correction}{0.42}           % AWWL's spectral correction, dB
\newcommand{\boostValue}{0.346}          % FID gained by applying it post hoc
\newcommand{\boostCI}{0.015}
\newcommand{\shortfall}{0.79}            % FID short of its own prediction
\newcommand{\fidMSE}{18.460}
\newcommand{\fidAWWL}{18.901}
\newcommand{\fidMSEboost}{18.114}
\newcommand{\fidAWWLboost}{18.568}

% Open questions, visible in the draft and impossible to forget at submission.
\newcommand{\todo}[1]{\textcolor{red}{\textbf{[TODO: #1]}}}
\newcommand{\pending}[1]{\textcolor{blue}{\textbf{[PENDING: #1]}}}

\title{Frequency-Aware Diffusion Losses Pay More Than They Buy}

\author{%
  \todo{author list; declare the CITA 2026 conference version here}%
}
\date{Draft --- \today}

\begin{document}
\maketitle

\begin{abstract}
Diffusion models under-produce high spatial frequencies, and a growing family
of training objectives sets out to correct this by weighting wavelet or Fourier
sub-bands of the loss. We audit one such objective over five seeds and find
that although it does exactly what it claims---closing \correction~dB of a
\deficitHigh~dB high-frequency deficit on CIFAR-10---it is no better than plain
MSE on FID, Inception Score or KID, and slightly worse on all three.

We then quantify what that correction should have been worth. Applying the
identical correction to the baseline's own samples as a Fourier post-process
improves FID by $\boostValue \pm \boostCI$, so the trained objective ends
\shortfall~FID short of what its own mechanism predicts: it buys a real
spectral improvement and pays for it several times over elsewhere. A twenty-line
post-process on the baseline outperforms the trained objective by that same
margin.

We further show that the published weighting conflates frequency balance with a
Min-SNR-style timestep reweighting ($|r| = 0.82$); isolating the frequency
component while holding the gradient budget fixed makes results significantly
worse ($+1.34$~FID, $p = 0.003$). Finally we report two measurement hazards that
reversed conclusions during this study: a hand-written KID estimator that
inverted the sign of the only positive result, and the observation that
calibrating metric sensitivity on real images measures it at a minimum where it
is quadratically insensitive, understating operating-point sensitivity roughly
thirtyfold.

\end{abstract}

\section{Introduction}
\label{sec:intro}

Samples from diffusion models are spectrally deficient, and measurably so. On
CIFAR-10 the radial power spectrum of DDPM samples sits \SI{0.36}{\decibel}
below that of real images at low frequency and \deficitHigh~\si{\decibel} below
at high frequency---a sixfold tilt with frequency
(\S\ref{sec:mechanism}). The observation is not new
\citep{durall2020watch,dzanic2020fourier,schwarz2021frequency}, and it has
motivated a family of training objectives that weight wavelet or Fourier
sub-bands of the loss, often on a schedule tied to the diffusion noise level
\citep{jiang2021focal,phung2023wavelet,huang2024wavedm}.

The premise is sound. This paper asks whether the remedy works, and---when it
does not---why, in terms precise enough to apply to the family rather than to
one implementation.

\paragraph{Contributions.}
\begin{enumerate}
  \item An analytic characterisation of the objective family: on an
    orthonormal basis, sub-band weighting reallocates a fixed error budget and
    adds no information (\S\ref{sec:family}).
  \item A confound in the standard weighting, which is simultaneously a
    timestep reweighting, and a protocol that separates the two
    (\S\ref{sec:confound}).
  \item A five-seed audit showing the mechanism works while every metric
    stays flat (\S\ref{sec:replication}, \S\ref{sec:mechanism}).
  \item A direct measurement of what the correction is worth, and the finding
    that the trained objective falls \shortfall~FID short of it
    (\S\ref{sec:worth}).
  \item Two measurement hazards that each reversed a conclusion mid-study and
    that generalise beyond this paper (\S\ref{sec:replication},
    \S\ref{sec:worth}).
\end{enumerate}

\paragraph{Relation to a prior conference version.}
The objective audited here was introduced in a conference paper by a subset of
the present authors \citep{thao2026awwl}, which reported gains on CIFAR-10 and
on subject-driven fine-tuning from single runs. This paper is not an extension
of that work in the usual sense: it re-evaluates those claims under seed
replication and does not reproduce them. We report the disagreement in full,
including two measurement errors of our own that had to be corrected before the
comparison could be trusted.
\todo{confirm the target venue's policy on extensions that revise a prior result}

\section{Related work}
\label{sec:related}

Four lines of work converge on the question this paper asks: whether a
measurable spectral correction translates into a measurable improvement.

\subsection{Spectral characterisation of generative models}

Generated images differ from real ones in their Fourier statistics, and the
difference is systematic. \citet{durall2020watch} traced a high-frequency
deficit in GAN samples to the up-convolution operators in the decoder;
\citet{dzanic2020fourier} showed the discrepancy is consistent enough across
architectures to serve as a detection signal, and \citet{chandrasegaran2021closer}
examined how much of it is intrinsic to the model rather than an artefact of
the last upsampling layer. \citet{wang2025spectral} supply the account of why
the deficit exists in the first place: solving the gradient-flow dynamics of a
denoiser yields an inverse-variance spectral law in which a Fourier mode's
learning time scales as $\tau \propto \lambda^{-1}$, so low-variance fine
detail is acquired orders of magnitude later than coarse structure. A
high-frequency deficit is therefore what an under-converged diffusion model
should look like, which is a competing explanation any frequency-aware
objective has to displace rather than assume. \citet{schwarz2021frequency} isolated generator and
discriminator separately and found the bias is not one effect: upsampling
operators bias the generator's spectrum, but checkerboard artefacts alone do
not explain the discrepancy, and the discriminator's difficulty is with
frequencies of \emph{low magnitude} rather than with high frequencies as such.

This literature establishes the phenomenon and, largely implicitly, a
motivation: because the spectrum is wrong, correcting it should help. That
implication is rarely tested directly, and it is what \S\ref{sec:worth} puts a
number on.

\subsection{Frequency-aware objectives and architectures}

Two families respond to the observation. The first modifies the
\emph{objective}. Focal Frequency Loss \citep{jiang2021focal} weights Fourier
components by how badly they are currently reconstructed, dynamically
emphasising the hard ones. Wavelet variants replace the Fourier basis with a
multi-resolution one, gaining spatial localisation: WaveDM \citep{huang2024wavedm}
applies wavelet supervision to restoration, and the objective audited here
weights wavelet sub-bands on a schedule tied to $\sigma_t$.

The second modifies the \emph{architecture} or the sampling path. WaveDiff
\citep{phung2023wavelet} moves the diffusion process itself into the wavelet
domain, primarily for speed; FreeU \citep{si2024freeu} rebalances the
backbone and skip contributions inside the U-Net at inference, and
\citet{yu2025frequency} rescale wavelet sub-bands during sampling to mitigate
exposure bias, and \citet{oh2026attention} edit cross-attention logits in the
Fourier domain at inference, one level deeper again. All three change the
frequency content of samples without touching training, which makes them the
closest relatives of the measurement in \S\ref{sec:worth} --- though there the
operation is used to price a mechanism rather than to improve one.
\citet{jiralerspong2025shaping} take a third route, shaping the frequency
content through the noising operator itself.

A recent comparative study of diffusion objectives \citep{kumar2025loss}
unifies several losses under the variational bound; frequency-weighted
objectives sit outside that unification, since their weights are not derived
from the bound but chosen to target a statistic of the output.

The family remains active. \citet{chandran2026spectral} add differentiable
Fourier \emph{and} wavelet losses to diffusion training as soft inductive
biases for frequency balance --- the closest contemporary sibling of the
objective we audit --- and report their largest gains on higher-resolution
unconditional data, where fine structure is hardest to model. That is worth
noting twice: it is independent support for the resolution caveat we state in
\S\ref{sec:limitations}, and it is a prediction the test of \S\ref{sec:worth}
could settle directly.

Reported gains in this family are typically small in absolute terms. The
improvements in \citet{chandran2026spectral} range from \num{0.02} to
\num{0.07} FID against baselines near \num{1.8}--\num{2.5}, averaged over three
seeds; those in the objective we audit are of comparable relative size. Our
results suggest that at that scale the question of whether an effect is present
at all requires the seed variance to be quantified first, and that even a
genuine spectral correction can be outweighed by what obtaining it costs.

\subsection{Timestep reweighting}

A parallel line reweights the loss across the noise schedule rather than across
frequency. \citet{nichol2021improved} resample timesteps by loss magnitude;
P2 \citep{choi2022perception} up-weights the region of the schedule where
perceptually relevant content is formed; Min-SNR \citep{hang2023efficient}
treats the schedule as a multi-task problem and clamps the per-timestep weight
by signal-to-noise ratio; \citet{karras2022elucidating} recast the whole design
space, weighting included, and \citet{kingma2021variational} derive a weighting
from the variational bound. \citet{li2026smoothed} unify the family from a
divergence-modification view and give an explicit cross-walk to Min-SNR, which
is what makes it meaningful to say a schedule \emph{is} a member of it rather
than merely correlated with one.

These two axes --- frequency and timestep --- are usually presented as
orthogonal. \S\ref{sec:confound} shows they are not orthogonal in the
formulation we audit: its frequency parameter simultaneously imposes a timestep
schedule correlated with Min-SNR at $|r| = 0.82$, so a result attributed to
frequency balance may belong to the other axis entirely. The multi-task framing
also supplies the natural alternative to hand-designed weights --- uncertainty
weighting \citep{kendall2018multi} and GradNorm \citep{chen2018gradnorm} --- which
\S\ref{sec:generality} evaluates.

\subsection{What the metrics can see}

A separate literature questions whether these metrics register what we assume.
\citet{barratt2018note} catalogued the Inception Score's failure modes;
\citet{chong2020effectively} showed FID is biased at finite sample size in a way
that depends on the model being evaluated; \citet{kynkaanniemi2023role}
demonstrated that FID responds strongly to the ImageNet class distribution of
the generated set, a property unrelated to image quality; and
\citet{stein2023exposing} found that the standard metrics rank diffusion models
against GANs in ways human raters do not reproduce.

Our contribution to this line is quantitative and specific: we measure how much
FID moves for a known spectral perturbation, and find that the answer depends
by more than an order of magnitude on whether the calibration is performed at
the metric's minimum or at a model's actual operating point
(\S\ref{sec:worth}).

\subsection{Replication and variance}

Finally, the practice this paper follows. \citet{bouthillier2021accounting}
showed that benchmark conclusions in machine learning frequently do not survive
proper accounting for sources of variance, seeds among them, and recommend
reporting comparisons with the variance made explicit.

\citet{dufour2026fid} put a number on it for exactly our setting. Across
several hundred networks trained on class-conditional ImageNet at
$256\times256$, retraining under an identical recipe with a different seed
moves FID $3.2\times$ more than redrawing samples from a fixed network, and
the coefficient of variation stays inside a \SIrange{1}{2}{\percent} band
regardless of model size or compute. They recommend treating any FID gap below
roughly \SI{1.3}{\percent} as inconclusive and reporting an error bar over
training seeds.

The agreement with our own measurement is worth stating because nothing about
it was arranged: our five-seed standard deviation of \num{0.21} at a baseline
FID of \fidMSE{} is a coefficient of variation of \SI{1.14}{\percent}, inside
their band, obtained on a different architecture, dataset and resolution. Two
independent estimates of how much a training seed is worth agreeing to within
a few tenths of a percent is the strongest evidence available that the quantity
is a property of diffusion training rather than of either setup. Differences of
the size routinely reported in the frequency-aware literature sit at or below
it, which is the methodological premise of \S\ref{sec:protocol}.

The 2026 sweep found no work that prices a training objective's stated
mechanism against a post-processed baseline, which remains this paper's
methodological contribution. Every cited entry has been verified against the
publisher or arXiv record.

\section{Background}
\label{sec:background}

\subsection{Diffusion training objectives}

A denoising diffusion model \citep{ho2020denoising} corrupts a datum $x_0$ by
$x_t = \sqrt{\bar\alpha_t}\,x_0 + \sqrt{1-\bar\alpha_t}\,\epsilon$ with
$\epsilon \sim \mathcal{N}(0, I)$, and trains a network
$\hat\epsilon_\theta(x_t, t)$ to recover $\epsilon$ under
\begin{equation}
  \mathcal{L}_{\mathrm{simple}}
  = \mathbb{E}_{x_0,\epsilon,t}\big[\|\hat\epsilon_\theta(x_t,t) - \epsilon\|_2^2\big].
  \label{eq:simple}
\end{equation}
We write $\sigma_t = \sqrt{1-\bar\alpha_t} \in (0,1)$ for the noise magnitude,
so $\sigma_t \to 1$ is the high-noise end of the schedule.

Equation~\ref{eq:simple} is a uniformly weighted regression: every timestep and
every spatial frequency contributes in proportion to its squared error. Both
uniformities have been questioned, and the two lines of work are usually
pursued separately --- a separation this paper argues is not maintained in
practice (\S\ref{sec:confound}).

\subsection{Wavelet sub-bands}

A one-level 2-D discrete wavelet transform \citep{mallat1989theory} splits an
image into an approximation band $\mathrm{LL}$ and three detail bands
$\mathrm{LH}$, $\mathrm{HL}$, $\mathrm{HH}$ carrying horizontal, vertical and
diagonal high-frequency content, each at half resolution. Iterating on
$\mathrm{LL}$ gives a multi-level decomposition.

Two properties matter here. The transform is \emph{linear}, so applying it to a
prediction and a target and differencing is identical to applying it to their
residual --- the decomposition acts on the error, not on image content. And for
an orthonormal family (Haar, Daubechies, Symlets) it is \emph{energy
preserving}, which is what makes Eq.~\ref{eq:parseval} exact and bounds what
any sub-band weighting can achieve.

\subsection{Evaluation}

Fréchet Inception Distance \citep{heusel2017gans} fits Gaussians to Inception-V3
features of real and generated sets and reports their Fréchet distance;
Kernel Inception Distance \citep{binkowski2018demystifying} replaces this with
an unbiased polynomial-kernel MMD in the same feature space; Inception Score
\citep{salimans2016improved} uses only the generated set. Precision and recall
\citep{kynkaanniemi2019improved,naeem2020reliable} separate fidelity from
coverage.

All of these are known to be sensitive to implementation details that are not
part of their definition: resizing and compression
\citep{parmar2022aliased}, finite-sample bias \citep{chong2020effectively},
and the particular ImageNet classes the feature extractor happens to encode
\citep{kynkaanniemi2023role}. \S\ref{sec:replication} adds one more to that
list, and \S\ref{sec:worth} shows that the operating point at which a metric's
sensitivity is measured changes the answer by more than an order of magnitude.

\section{The objective family}
\label{sec:family}

Write $\hat{\epsilon}$ for the predicted noise and $\epsilon$ for the target.
A sub-band objective applies a linear transform $W$---here a 2-D discrete
wavelet transform---to both and weights the resulting bands:
\begin{equation}
  \mathcal{L} = \sum_{b} w_b(\sigma_t) \,\big\| (W\hat{\epsilon})_b - (W\epsilon)_b \big\|_2^2 .
  \label{eq:family}
\end{equation}

Because $W$ is linear, $W\hat{\epsilon} - W\epsilon = W(\hat{\epsilon} -
\epsilon)$: the bands being weighted are bands of the \emph{prediction
residual}, not of any image. This is worth stating plainly, because motivations
phrased in terms of separating image structure from texture describe something
the objective does not do---the $\epsilon$-prediction target is white noise and
has no structure to decompose.

For an orthonormal basis, Parseval's identity gives
\begin{equation}
  \tfrac{1}{4}\textstyle\sum_b \operatorname{mean}\big((W r)_b^2\big)
  = \operatorname{mean}(r^2), \qquad r = \hat{\epsilon} - \epsilon,
  \label{eq:parseval}
\end{equation}
verified numerically to a relative error of $7.4\times 10^{-8}$ for the Haar
basis. An \emph{unweighted} sub-band loss is therefore exactly the pixel MSE,
and every member of the family redistributes one fixed error budget across
orthogonal bands. Any gain must come from \emph{how} the budget is split; none
can come from added information.

\section{A confound, and a protocol}
\label{sec:confound}

The weighting we audit takes the form
\begin{equation}
  w_{\mathrm{LL}}(\sigma) = \frac{\alpha\,\sigma^{p}}{\sigma^{p} + (1-\sigma)^{p}},
  \qquad
  w_{\mathrm{det}}(\sigma) = \frac{(1-\alpha)(1-\sigma)^{p}}{\sigma^{p} + (1-\sigma)^{p}} .
  \label{eq:weights}
\end{equation}

The two share a denominator but do \emph{not} sum to a constant: the total runs
from $\alpha$ as $\sigma \to 1$ to $1-\alpha$ as $\sigma \to 0$. So $\alpha$
sets the frequency balance \emph{and} rescales gradient magnitude across the
noise schedule. Measured against the Min-SNR-$\gamma$ curve
\citep{hang2023efficient} over the $\sigma_t$ of all $1000$ training timesteps,
the correlation is $r = +0.82$ at $\alpha = 0.8$ and $r = -0.82$ at
$\alpha = 0.2$: the two reported optima are opposite timestep schedules, not
merely different frequency balances (Fig.~\ref{fig:weights}). The figure is
sensitive to the grid --- a
uniform grid in $\sigma$ gives $0.87$ --- so we quote it on the distribution
training actually samples.

\guardedfigure{0.95}{weights.png}{%
  The published weighting at $\alpha = 0.2$, $p = 1$. The two band weights
  (blue, orange) are what the objective is described as doing; their total
  (black) is what it also does, and it is not constant. Plotted against the
  Min-SNR-$\gamma$ curve (violet) with both normalised to unit mean, the total
  is visibly a timestep schedule of the same family. Sampled at the $\sigma_t$
  of the \num{1000} training timesteps.}{fig:weights}

\paragraph{Protocol.} To vary only the shape, hold $\sum_b w_b$ constant
\emph{at its original mean} of $0.5$. Normalising naively to $1$ doubles the
mean gradient magnitude and confounds the ablation with a $2\times$ learning
rate --- a mistake we made first and caught by measuring.

\begin{table}[t]
\centering
\caption{Isolating the frequency component makes matters worse. CIFAR-10,
five seeds, epoch 199; $p$ is Holm-corrected across the comparison family.}
\label{tab:confound}
\begin{tabular}{lrr}
\toprule
Configuration & $\Delta$ FID vs.\ MSE & $p$ (Holm) \\
\midrule
Published weighting & $+0.44$ & $0.43$ \\
Normalised to $1$ (confounded, $2\times$ LR) & $+1.15$ & $0.045$ \\
\textbf{Budget held fixed} & $\mathbf{+1.34}$ & $\mathbf{0.003}$ \\
\bottomrule
\end{tabular}
\end{table}

Table~\ref{tab:confound} settles the question: separated from the timestep
effect, the frequency weighting is significantly \emph{harmful}. What residual
appeal the published form has comes from the timestep reweighting riding along
with it.

\section{Experimental protocol}
\label{sec:protocol}

Every configuration is trained with five seeds; we report mean $\pm$ standard
deviation with a 95\% confidence interval, and compare with a paired-by-seed
$t$-test under Holm--Bonferroni correction across the comparison family.
Pairing removes the variance the seed itself contributes and is markedly more
sensitive at $N=5$ than an unpaired test.

We also report a Wilcoxon signed-rank test, but as a direction check only: its
smallest attainable two-sided $p$ is $2/2^{N} = 0.0625$ at $N=5$, so it cannot
clear $\alpha = 0.05$ and must not be read as a verdict.

Sampling uses DDIM \citep{song2021denoising} with 100 steps and 10\,000 samples,
identical across arms; FID via \texttt{clean-fid} \citep{parmar2022aliased}, IS
and KID via \texttt{torch-fidelity}. Spectra are reported in decibels
($20\log_{10}|F|$). Exponential moving averaging is reported as its own
configuration rather than mixed in.

\pending{DDPM-1000 sampler sensitivity check --- the one open threat to the null}

\section{The metrics do not move}
\label{sec:replication}

\begin{table}[t]
\centering
\caption{CIFAR-10, 200 epochs, five seeds. No metric favours the
frequency-aware objective; all three point the other way.}
\label{tab:replication}
\begin{tabular}{lrr}
\toprule
Metric & $\Delta$ vs.\ MSE & $p$ (Holm) \\
\midrule
FID $\downarrow$ & $+0.44$ (worse) & $0.43$ \\
IS $\uparrow$ & $-0.059$ (worse) & $0.44$ \\
KID $\downarrow$ & $+0.0003$ (worse) & $0.77$ \\
\bottomrule
\end{tabular}
\end{table}

Table~\ref{tab:replication} reports the audit. Two details deserve emphasis.

First, the conference version reports an Inception Score gain of $+0.17$. Our
95\% confidence interval on that difference is $[-0.171, +0.054]$, which
\emph{excludes} the reported value. The data do not merely fail to confirm the
gain; they are inconsistent with it.

Second, convergence is \emph{slower} at every checkpoint (FID $59.9$ vs.\
$52.4$ at epoch 39), which closes the natural fallback claim of reaching a
given quality with less compute.

\paragraph{Equation versus code.} Implementing Eq.~\ref{eq:family} exactly as
the conference version writes it --- summing the three detail bands rather than
averaging them --- costs $+2.63$~FID ($p = 0.0004$). The published equation and
the released implementation are different objectives, and the equation is much
the worse of the two. A reader reimplementing from the paper obtains something
worse than the baseline.

\paragraph{Measurement hazard: a non-standard estimator reversed the sign.}
A hand-written polynomial-kernel KID showed the objective \emph{better}
($p = 0.028$, uncorrected), and it was the only positive result in the study.
The same quantity computed by an established library shows it \emph{worse}.
Because both FID and KID measure distributional distance in the same Inception
feature space, the disagreement was itself the signal that something was wrong.
We report the library value throughout, and recommend that any headline claim
resting on a bespoke estimator be cross-checked before it is believed.

\section{Results}
\label{sec:results}

All tables in this section are generated directly from the results ledger by
\texttt{scripts/make\_tables.py}; none of the numbers are transcribed. Arms
averaged over fewer than five seeds carry a superscript.

% Generated tables. \IfFileExists keeps the draft building on a machine that
% has the manuscript but not the sweep output.
\IfFileExists{tables/full.tex}{% Generated by scripts/make_tables.py --- do not edit by hand.
% Regenerate after every sweep; the paper reads these directly.

\begin{table}[t]
\centering
\small
\caption{CIFAR-10, 200 epochs, epoch 199. Mean $\pm$ standard deviation over 5 seeds; a superscript marks arms with fewer. KID is from \texttt{torch-fidelity}; see \S\ref{sec:replication} for why the bespoke estimator is not reported.}
\label{tab:full}
\begin{tabular}{lrrrrr}
\toprule
Objective & FID $\downarrow$ & IS $\uparrow$ & KID $\downarrow$ ($\times 10^{3}$) & Prec. $\uparrow$ & Rec. $\uparrow$ \\
\midrule
MSE \emph{(baseline)} & $18.46 \pm 0.21$ & $7.735 \pm 0.070$ & $13.34 \pm 0.27$ & $0.776 \pm 0.002$ & $0.641 \pm 0.003$ \\
$L_1$ & $25.60 \pm 0.17$ & $7.681 \pm 0.057$ & $18.01 \pm 0.20$ & $0.760 \pm 0.003$ & $0.608 \pm 0.004$ \\
Huber & $18.80 \pm 0.02$\,\textsuperscript{2} & $7.769 \pm 0.059$\,\textsuperscript{2} & $13.90 \pm 0.08$\,\textsuperscript{2} & $0.775 \pm 0.001$\,\textsuperscript{2} & $0.647 \pm 0.002$\,\textsuperscript{2} \\
Static wavelet ($p{=}0$) & $18.63 \pm 0.06$ & $7.748 \pm 0.019$ & $13.64 \pm 0.16$ & $0.780 \pm 0.005$ & $0.641 \pm 0.004$ \\
Frequency-aware \emph{(published)} & $18.90 \pm 0.34$ & $7.676 \pm 0.060$ & $13.65 \pm 0.45$ & $0.779 \pm 0.003$ & $0.634 \pm 0.003$ \\
\quad as Eq.~(1) is written & $21.09 \pm 0.29$ & $7.500 \pm 0.023$ & $16.06 \pm 0.24$ & $0.782 \pm 0.006$ & $0.623 \pm 0.004$ \\
\quad normalised to 1 & $19.61 \pm 0.40$ & $7.578 \pm 0.035$ & $14.64 \pm 0.32$ & $0.779 \pm 0.005$ & $0.631 \pm 0.004$ \\
\quad budget held fixed & $19.80 \pm 0.19$ & $7.595 \pm 0.053$ & $14.80 \pm 0.17$ & $0.776 \pm 0.004$ & $0.630 \pm 0.004$ \\
MSE + EMA & $18.46 \pm 0.37$\,\textsuperscript{3} & $7.724 \pm 0.079$\,\textsuperscript{3} & $13.46 \pm 0.25$\,\textsuperscript{3} & $0.778 \pm 0.001$\,\textsuperscript{3} & $0.642 \pm 0.001$\,\textsuperscript{3} \\
Frequency-aware + EMA & $18.95 \pm 0.20$\,\textsuperscript{3} & $7.683 \pm 0.065$\,\textsuperscript{3} & $13.74 \pm 0.34$\,\textsuperscript{3} & $0.776 \pm 0.002$\,\textsuperscript{3} & $0.637 \pm 0.001$\,\textsuperscript{3} \\
\bottomrule
\end{tabular}
\end{table}
}{%
  \begin{table}[t]\centering
  \fbox{\parbox{0.8\linewidth}{\centering\vspace{4mm}
  \textbf{tables/full.tex not built}\\[1mm]
  run \texttt{python scripts/make\_tables.py}\vspace{4mm}}}
  \caption{Full results across objectives and metrics.}\label{tab:full}
  \end{table}}

% The fallback carries the label as well as the box: without it every \ref to
% the table dangles, and a dangling reference is easy to stop noticing.
\IfFileExists{tables/compare_fid.tex}{\input{tables/compare_fid}}{%
  \begin{table}[t]\centering
  \fbox{\parbox{0.8\linewidth}{\centering\vspace{4mm}
  \textbf{tables/compare\_fid.tex not built}\\[1mm]
  run \texttt{python scripts/make\_tables.py}\vspace{4mm}}}
  \caption{Paired comparison against MSE on FID.}\label{tab:compare-fid}
  \end{table}}

Table~\ref{tab:full} reports every objective on every metric;
Table~\ref{tab:compare-fid} gives the paired comparison against MSE. Three
readings matter.

\paragraph{The published objective is not better on any metric.}
Its FID, Inception Score and KID all sit on the wrong side of the baseline, and
none of the three differences clears significance. The point estimates are
consistent with each other in sign, which is worth more than any of them
individually: three metrics computed from two different feature pipelines agree
on the direction.

\paragraph{The ablations that were supposed to isolate the mechanism make it
worse, significantly.} Both normalised variants and the literal reading of the
published equation are significantly worse than MSE after correction, while the
published form itself is not distinguishable from MSE. The ordering is
informative: the closer the objective gets to being \emph{purely} a frequency
weighting, the worse it performs. What keeps the published form competitive is
the timestep reweighting it performs incidentally (\S\ref{sec:confound}).

\paragraph{Seed variance dominates the differences the literature reports.}
The standard deviation of FID across five seeds of an identical configuration
is \num{0.21} for the baseline. Improvements reported in this area are
routinely smaller than that. Any single-run comparison at this scale is
uninformative regardless of which way it happens to fall, which is the
methodological point of \citet{bouthillier2021accounting} instantiated for
diffusion training.

\subsection{Convergence}

The natural fallback for an objective that does not improve final quality is
that it reaches a given quality sooner. It does not: FID is worse at every
evaluated checkpoint (\num{59.9} against \num{52.4} at epoch~39, \num{36.9}
against \num{33.9} at epoch~79, and so on to epoch~199). The gap narrows as
training proceeds but never closes or reverses.

\subsection{Exponential moving averaging}

EMA is reported as its own configuration because it is standard practice and
the published runs omit it, so mixing the two would confound the comparison.
It shifts both arms by less than the seed variance and leaves their ordering
unchanged, which is expected: at a decay of \num{0.9999} over
\num{78200} steps the shadow copy averages roughly the last \num{10000} steps,
by which point the iterate is moving slowly.

\IfFileExists{tables/compare_is.tex}{% Generated by scripts/make_tables.py --- do not edit by hand.
% Regenerate after every sweep; the paper reads these directly.

\begin{table}[t]
\centering
\small
\caption{Paired-by-seed comparison against MSE on Inception Score. $p$ is Holm--Bonferroni corrected across the rows shown. Wilcoxon cannot fall below $0.0625$ at five seeds and is a direction check, not a verdict.}
\label{tab:compare-ismean}
\begin{tabular}{lrrrl}
\toprule
Objective & $\Delta$ & $t$ & $p$ (Holm) & Verdict \\
\midrule
\quad normalised to 1 & $-0.157$ & $-9.98$ & $0.0051$ & \textbf{worse} \\
\quad as Eq.~(1) is written & $-0.235$ & $-7.25$ & $0.0154$ & \textbf{worse} \\
\quad budget held fixed & $-0.140$ & $-5.66$ & $0.0336$ & \textbf{worse} \\
Frequency-aware \emph{(published)} & $-0.059$ & $-1.45$ & $1.0000$ & n.s. \\
Frequency-aware + EMA & $-0.041$ & $-0.62$ & $1.0000$ & n.s. \\
Huber & $-0.009$ & $-0.19$ & $1.0000$ & n.s. \\
$L_1$ & $-0.054$ & $-1.18$ & $1.0000$ & n.s. \\
MSE + EMA & $+0.001$ & $0.08$ & $1.0000$ & n.s. \\
Static wavelet ($p{=}0$) & $+0.013$ & $0.36$ & $1.0000$ & n.s. \\
\bottomrule
\end{tabular}
\end{table}
}{%
  \begin{table}[t]\centering
  \fbox{\parbox{0.8\linewidth}{\centering\vspace{4mm}
  \textbf{tables/compare\_is.tex not built}\\[1mm]
  run \texttt{python scripts/make\_tables.py}\vspace{4mm}}}
  \caption{Paired comparison against MSE on Inception Score.}
  \label{tab:compare-ismean}
  \end{table}}

\section{The mechanism works}
\label{sec:mechanism}

\begin{table}[t]
\centering
\caption{Signed deviation from the real radial spectrum in decibels; negative
means too little energy. The static ablation is indistinguishable from MSE.}
\label{tab:bands}
\begin{tabular}{lrrr}
\toprule
Configuration & Low & Mid & High \\
\midrule
MSE & $-0.36$ & $-0.80$ & $-2.20$ \\
Static wavelet ($p{=}0$) & $-0.35$ & $-0.79$ & $-2.23$ \\
Frequency-aware & $-0.31$ & $-0.56$ & $\mathbf{-1.79}$ \\
\bottomrule
\end{tabular}
\end{table}

\begin{figure}[t]
\centering
\includegraphics[width=0.85\linewidth]{spectrum.png}
\caption{Deviation from the real spectrum against normalised frequency. The
MSE and static-wavelet curves coincide across the entire axis: the wavelet
decomposition on its own changes nothing. Only the time-dependent schedule
separates from them, and the gap widens with frequency.}
\label{fig:spectrum}
\end{figure}

Table~\ref{tab:bands} and Figure~\ref{fig:spectrum} establish three things.

The premise holds: the deficit is roughly sixfold worse at high frequency than
at low. The objective's improvement is largest there in absolute terms
(\SI{0.41}{\decibel} against \SI{0.05}{\decibel} at low frequency), so the
correction lands where the method says it should. We state this precisely: the
improvement is \emph{broadband with a high-frequency tilt}, not an exclusively
high-frequency correction --- the proportional gain is largest in the mid band,
and a referee will recompute.

The wavelet decomposition contributes nothing. The static ablation, which keeps
the decomposition and removes only the dependence on $\sigma$, overlaps MSE
across the whole frequency axis. The temporal schedule is the entire effect.
The correction also survives the protocol of \S\ref{sec:confound}
($p = 0.0056$), so it is a genuine frequency effect rather than the timestep
confound in disguise.

\begin{figure}[t]
\centering
\includegraphics[width=\linewidth]{compare.png}
\caption{Samples from identical initial noise, one row per objective. A
\SI{0.4}{\decibel} high-frequency change at $32\times32$ is below the visible
threshold.}
\label{fig:compare}
\end{figure}

Finally, the correction is not visible. Figure~\ref{fig:compare} pairs samples
on their initial latent, so any difference would be attributable to the
objective; none is apparent.

\section{What the correction is worth}
\label{sec:worth}

Sections \ref{sec:replication} and \ref{sec:mechanism} leave an obvious
question unanswered: how much \emph{should} FID have moved? Without that
number, ``the metrics did not respond'' is an observation rather than an
argument.

We answer it by applying the correction directly. Taking each model's own
samples and boosting the high band by \SI{0.41}{\decibel} --- the amount the
trained objective achieves --- is a pure post-process, so the resulting change
in FID isolates the value of the spectral correction from everything else
training might have altered.

\begin{table}[t]
\centering
\caption{The correction applied as a Fourier post-process, five seeds per arm.
The \SI{0}{\decibel} column reproduces the sweep's FID exactly, validating the
apparatus.}
\label{tab:worth}
\begin{tabular}{lrrr}
\toprule
& Spectral deficit & FID & FID after $+\SI{0.41}{\decibel}$ \\
\midrule
MSE & \SI{2.22}{\decibel} & $\fidMSE \pm 0.210$ & $\mathbf{\fidMSEboost \pm 0.202}$ \\
Frequency-aware & \SI{1.80}{\decibel} & $\fidAWWL \pm 0.344$ & $\fidAWWLboost \pm 0.335$ \\
\bottomrule
\end{tabular}
\end{table}

Two independent measurements agree that the objective corrects
$\correction \pm 0.18$~\si{\decibel}: the difference of deficits in
Table~\ref{tab:worth}, and the band analysis of \S\ref{sec:mechanism} obtained
by a different route.

Applying that correction post hoc is worth $\boostValue \pm \boostCI$~FID. The
accounting follows immediately.

\begin{table}[t]
\centering
\caption{The objective falls \shortfall~FID short of what its own spectral
correction predicts.}
\label{tab:accounting}
\begin{tabular}{lr}
\toprule
& FID \\
\midrule
MSE baseline & \fidMSE \\
Predicted for the objective from its own spectral correction & \fidMSEboost \\
Objective, measured & \fidAWWL \\
\midrule
\textbf{Unexplained penalty} & $\mathbf{+\shortfall}$ \\
\bottomrule
\end{tabular}
\end{table}

The objective buys $0.35$~FID of spectral improvement and pays roughly
$1.1$~FID for it elsewhere. Two framings of the same fact are worth stating
separately:

\begin{itemize}
  \item MSE plus a twenty-line Fourier post-process (\fidMSEboost) outperforms
    the trained objective (\fidAWWL) by \shortfall~FID, at no training cost.
  \item The objective plus the same post-process (\fidAWWLboost) is
    \emph{still} worse than plain MSE (\fidMSE). The correction cannot recover
    what obtaining it costs.
\end{itemize}

\paragraph{Measurement hazard: calibrate away from the minimum.}
We first calibrated FID's sensitivity on real images --- attenuating one half
of a real set and scoring it against the other. That measures sensitivity at
the \emph{minimum} of the distance, where the first derivative vanishes and
perturbations are penalised quadratically. It implied a \SI{0.41}{\decibel}
correction was worth $0.010$~FID and would require some 17\,500 seeds to
detect. Measured at the operating point, where the derivative is non-zero, the
same correction is worth $\boostValue$ --- thirtyfold larger. Sensitivity
calibrations for a generative metric must be performed away from the minimum;
we drew, and retract, the opposite conclusion from the wrong calibration.

\section{Generality}
\label{sec:generality}

\pending{tier-4 learned weightings on CIFAR-10: uncertainty-conditioned and
GradNorm}
\pending{second resolution (CelebA-64) --- required for a top-tier claim}
\pending{second architecture (latent diffusion or DiT) --- required likewise}

Everything so far concerns one task: unconditional generation trained from
scratch at $32\times32$. The objective's published optimum on a second task ---
subject-driven fine-tuning with DreamBooth \citep{ruiz2023dreambooth} --- is
$\alpha = 0.8$, the mirror image of the $\alpha = 0.2$ used above, and that
inversion is the evidence offered for a resolution-dependent law. It is also
the setting where \S\ref{sec:confound} predicts something specific: if
$\alpha$'s effect runs through the timestep axis rather than the frequency
one, then a pure timestep method should do at least as well as the
frequency-aware objective here, and the frequency-aware objective's advantage
should not survive replication.

\subsection{The published table, replicated}

We fine-tune Stable Diffusion~1.5 on one subject for \num{400} steps, five
seeds per objective, and score \num{60} generated images per run against three
prompts: CLIP text-image agreement, and CLIP image-image similarity to the
instance set the run was trained on.

\IfFileExists{tables/dreambooth.tex}{\input{tables/dreambooth}}{%
  \begin{table}[t]\centering
  \fbox{\parbox{0.8\linewidth}{\centering\vspace{4mm}
  \textbf{tables/dreambooth.tex not built}\\[1mm]
  run \texttt{python scripts/make\_tables.py}\vspace{4mm}}}
  \caption{DreamBooth, five seeds per objective.}\label{tab:dreambooth}
  \end{table}}

The published claim is that the audited objective is the only method in the
top two on \emph{both} metrics. Over five seeds it is not: it ranks first of
nine on subject similarity (\num{0.8721}) and seventh on prompt agreement
(\num{0.3003}), below the MSE baseline. Neither of its differences from MSE
survives correction --- $-0.0108$ on CLIP score and $+0.0292$ on similarity,
the latter at $p = 0.050$ uncorrected and $p = 0.35$ after Holm. As on
CIFAR-10, the reported ordering is a reading of differences smaller than the
seed variance that produces them.

\subsection{The timestep axis is the one that moves}

One row does clear correction, and it is the row that carries no frequency
component at all. Min-SNR-$\gamma$ weighting \citep{hang2023efficient}
improves CLIP score by $+0.0411$ ($p = 0.005$, Holm), four times the largest
effect any wavelet variant produces and in the opposite direction to most of
them.

This is the prediction of \S\ref{sec:confound} confirmed on an independent
task. There, the published $\alpha$ was shown to impose a timestep schedule
correlated with Min-SNR at $|r| = 0.82$, so that a result credited to
frequency balance might belong to the other axis. Here the two axes are
separated by construction --- one objective reweights only frequency, another
only timesteps --- and only the timestep one produces an effect the data can
distinguish from noise.

The effect is a trade rather than an improvement. Min-SNR is simultaneously
the best objective on prompt agreement and the worst on subject similarity
($-0.0441$ against MSE), which is what a shift in where gradient mass falls
along the schedule should look like: prompt-following and subject fidelity are
formed at different noise levels, and emphasising one costs the other. We
report it as a relocation of the trade-off, not as a better objective.

\subsection{What five seeds can and cannot establish}

The paired $t$-test reports $p = 0.0006$ for the Min-SNR effect on CLIP score.
The Wilcoxon signed-rank test on the same five pairs reports $p = 0.0625$,
which is its \emph{smallest attainable value} at $n = 5$: the statistic is
bounded below by $2/2^{n}$, so no rank-based test on five seeds can reach
\num{0.05} however large the effect. The two numbers do not disagree; the
second is at its floor.

We report both throughout for that reason. A rank test is the natural
robustness check against the $t$-test's normality assumption, and at the seed
counts common in this literature it cannot supply one. Reading its $p$-value
as evidence of absence --- rather than as the floor it sits on --- would
reverse this section's only positive finding.

\section{Discussion}
\label{sec:discussion}

\subsection{A mechanism can be real and still not pay}

The result that organises this paper is not that a frequency-aware objective
failed. It is that the objective \emph{succeeded} at the thing it was designed
to do and lost anyway. It closes \correction~\si{\decibel} of a
\deficitHigh~\si{\decibel} high-frequency deficit; the correction is
attributable to the component its design credits, since removing the temporal
schedule removes the whole effect (\S\ref{sec:mechanism}); and the correction
is worth $\boostValue$~FID when applied in isolation (\S\ref{sec:worth}). The
objective nonetheless lands \shortfall~FID short of what that predicts.

This is a different failure from the one usually reported. A method that does
not work is uninteresting; a method that does exactly what it claims and is
still outperformed by its baseline says something about the relationship
between the mechanism and the goal. Here the relationship is that the mechanism
is not free: reallocating the error budget towards high frequencies means
allocating less of it everywhere else, and by Eq.~\ref{eq:parseval} that is
precisely what the reallocation must do. What the objective gives up is not
visible in the spectrum, because the spectrum is the one thing it was
optimising.

\subsection{Post-processing as a baseline}

That the same spectral correction can be obtained by a Fourier operation on
finished samples --- outperforming the trained objective --- is not offered as
a method. It is offered as a control that this literature does not use and, we
think, should.

Any objective whose stated mechanism is a change to an image statistic admits
the same test: apply that change to the baseline's outputs directly and see
what it is worth. If the trained method does not clear the post-processed
baseline, the training-time machinery is not paying for itself, whatever the
ablations show. The test is cheap, it requires no additional training, and it
would have flagged the objective audited here before any of the sweeps in
\S\ref{sec:results} were run.

\subsection{What the metrics can and cannot register}

Our two measurement hazards share a structure: in both, a quantity was
estimated correctly under one set of conditions and read as though it held
under another.

The KID reversal (\S\ref{sec:replication}) is the familiar version --- a
bespoke estimator, correct in its own terms, disagreeing in sign with the
standard one. It is worth reporting because the disagreement was the study's
only positive result, and because the failure is invisible unless a second
implementation is run.

The calibration error (\S\ref{sec:worth}) is subtler and, we suspect, more
common. Measuring how much FID responds to a perturbation of real images
measures the response \emph{at the minimum}, where the first-order term
vanishes by construction and the metric is quadratically insensitive. It is a
natural experiment to design and it gives an answer roughly thirty times too
small. Sensitivity has to be probed where the model actually sits. We report
this because we drew the wrong conclusion from it first, and because the same
design suggests itself to anyone asking the same question.

Neither hazard is an argument against these metrics. Both are arguments for
asking what a result \emph{should} have been before deciding what it means.

\subsection{Implications for frequency-aware training}

We audit one objective, and we are careful not to generalise the empirical
result past it. But three parts of the analysis are not specific to it.

Equation~\ref{eq:parseval} constrains the whole family: on an orthonormal
basis, sub-band weighting redistributes a fixed budget, so every gain is
someone else's loss and the question is always whether the trade is favourable.
The confound of \S\ref{sec:confound} arises whenever a weighting depends on
$\sigma_t$ without its total being held fixed, which is the common case; the
protocol that separates the axes costs nothing to adopt. And the
post-processing control of \S\ref{sec:worth} applies to any method that names
the statistic it intends to change.

Whether other members of the family clear these bars is an empirical question
we have not answered. We would expect the answer to depend on resolution: at
$32\times32$ the high-frequency band carries little of the signal that
Inception features encode, and the trade a frequency-aware objective makes may
look different where it carries more.

\section{Limitations}
\label{sec:limitations}

\begin{itemize}
  \item A single architecture and dataset until \S\ref{sec:generality} lands.
    Absence of an effect at $32\times32$ does not establish absence at
    resolutions where high-frequency content carries more of the signal.
  \item Sampling is DDIM-100 throughout. The comparison is internally
    consistent, but a sampler with fewer steps produces smoother output, which
    is precisely where a frequency-aware objective should have its advantage.
    \pending{DDPM-1000 check}
  \item The spectral distance is computed by a hand-written radial FFT. Two
    independent implementations here agree, which is reassuring and not proof
    ---and after \S\ref{sec:replication} we are in no position to treat our own
    implementations as authoritative.
  \item The Fourier post-process of \S\ref{sec:worth} leaves a spatial
    signature unlike a model's under-production, so the value it assigns to the
    correction is an estimate rather than an exact accounting.
\end{itemize}

\section{Conclusion}
\label{sec:conclusion}

A frequency-aware diffusion objective can do exactly what it was designed to
do---measurably correcting a real spectral deficiency, through the component
its design credits---and still lose to the baseline it was meant to improve
on. The correction is worth something, and we measure what: $\boostValue$~FID.
The objective delivers it and gives back three times as much elsewhere, so that
a post-process obtaining the same correction for free outperforms it.

The broader lesson is about attribution. It is not enough to show that a method
changes the quantity it targets; the change must be priced against what the
evaluation metric can register, and against what the method costs to obtain it.
Both of the measurement hazards we report were caught only because we asked
what a result \emph{should} have been before accepting what it was.

\todo{revisit once the learned-weighting arms report}


\bibliographystyle{plainnat}
\bibliography{refs}

\end{document}
